\pagenumbering{arabic}
\section{引言}

\subsection{研究背景}

生物医学数据科学的核心特征在于，其问题求解依赖于可执行的计算流程，而非单纯的理论陈述。典型的数据科学任务涵盖了从数据清洗、结构化处理，到统计建模、结果可视化及验证的全过程。随着大语言模型（Large Language Model, LLM）在代码生成领域的快速发展，能够根据任务需求自动生成、执行并调试代码的智能系统，已成为处理真实生物医学数据流的关键工具。

然而，生物医学数据分析任务具有高度的领域特殊性：临床数据包含复杂的患者标识体系、异构的数据类型（基因组变异、生存数据、临床特征等），以及严格的统计验证要求。这些特点使得通用的代码生成方法难以直接适用，需要构建具备领域感知能力的智能体系统。

\subsection{任务建模}

本项目将生物医学代码推理任务建模为部分可观测马尔可夫决策过程（POMDP）。给定问题描述 $x \in \mathcal{X}$，智能体的目标是生成能够产出正确结果 $y = y^*$ 的代码 $c \in \mathcal{C}$，其中 $y^*$ 是标准答案。其交互逻辑遵循 ReAct（Reasoning and Acting）范式\cite{yao2023react}，通过"推理--行动--观测"的迭代循环展开：

\begin{equation}
(\text{Thought}_t \rightarrow a_{t+1} \rightarrow o_{t+1})_{t=0}^{T-1}
\end{equation}

其中 $\text{Thought}_t$ 是智能体在时间步 $t$ 基于观察 $o_t$ 进行的内部推理与策略规划；$a_{t+1}$ 指从预定义的原子动作空间 $\mathcal{A}$ 中采样的具体可执行动作；$o_{t+1}$ 是环境反馈的新观测值。最终，系统通过验证函数 $E(c, y) = \mathbb{I}(y = y^*)$ 判定执行输出与标准答案的一致性。

\subsection{主要贡献}

本项目的主要贡献如下：

\begin{enumerate}
    \item 实现了基于 ReAct 范式的生物医学编码智能体，包含完整的 Thought--Action--Observation 闭环控制逻辑，支持四类原子动作，并通过子进程隔离的安全沙箱执行代码，具备超时控制和危险操作拦截能力。
    \item 设计并实现了反思与经验复用池（Reflexion \& Episodic Memory）机制，通过结构化反思提取成功策略和失败教训，在新任务中检索相似经验作为 few-shot 注入，实现跨任务的长期记忆。
    \item 构建了多智能体协作系统，包含数据分析、建模、SQL 查询、质量保证四类专门化智能体，支持顺序、并行、分层、自适应四种协作模式。
    \item 开发了基于 Streamlit 的可视化 Web 监控面板，支持实时查看智能体执行轨迹和任务状态。
    \item 在 BioDSBench\cite{wang2024biodsbench} 数据集上进行了系统的消融实验，验证了经验池机制的有效性和模型能力对任务通过率的影响。
\end{enumerate}
